% ============================================================
% PAGE 3 — Related Work: Deep Learning Detectors
% ============================================================

\subsection{Two-Stage Deep Learning Detectors}

\subsubsection{R-CNN Family}

Girshick \textit{et al.}~\cite{girshick2014rich} introduced
Region-based CNN (R-CNN), the first detector to couple CNN
features with region proposals.  Selective search generates
approximately 2\,000 candidate regions per image; each region
is warped to a fixed $227 \times 227$ pixel window and
forwarded through an AlexNet-derived CNN to produce a 4\,096-dimensional
feature vector, which a class-specific SVM then classifies.
Although R-CNN substantially outperformed DPM on PASCAL VOC,
its inference pipeline was glacially slow: approximately
47 seconds per image on a GPU, because each of the 2\,000
proposals required a separate CNN forward pass.

Fast R-CNN~\cite{girshick2015fast} resolved the redundant
computation by feeding the full image through the CNN once to
produce a shared convolutional feature map, then projecting
each proposal onto this map via the RoI Pooling operation.  A
lightweight per-RoI network head simultaneously produces class
scores and bounding-box regression offsets for all proposals in
a single forward pass.  By replacing the independent per-patch
CNNs with shared feature extraction, Fast R-CNN reduced
inference time to $0.3$ seconds per image while improving
accuracy.

Faster R-CNN~\cite{ren2015faster} eliminated the external
selective-search proposal algorithm by inserting a Region
Proposal Network (RPN) that shares the convolutional backbone
with the detection head.  The RPN slides a small network over
the shared feature map, predicting, at each location, $k$
anchor-box objectness scores and associated regression offsets.
Alternating training of the RPN and the detection head allowed
end-to-end optimization.  Faster R-CNN achieved 5--17 FPS
at 73.2\% mAP on PASCAL VOC 2007---a milestone demonstrating
that deep detection could be practical rather than merely
accurate.  Subsequent variants---R-FCN, FPN-based Faster
R-CNN, Cascade R-CNN---progressively improved accuracy to
over 40 mAP on COCO by adding more sophisticated region
classifiers and multi-scale feature fusion, but FPS remained
limited by the inherently sequential two-stage computation.

\subsection{Single-Stage Deep Learning Detectors}

\subsubsection{SSD}

Liu \textit{et al.}~\cite{liu2016ssd} proposed Single Shot
MultiBox Detector (SSD), which eliminated the proposal stage
entirely.  SSD places a set of default anchor boxes at every
location of several feature maps extracted at different
resolutions of the backbone; for each anchor, a small
convolutional network predicts class probabilities and
box offsets.  By producing all predictions in a single forward
pass---without the sequential region-wise processing of
two-stage methods---SSD achieved 59 FPS at 74.3 mAP on PASCAL
VOC, demonstrating that real-time accuracy was attainable
without a proposal network.  Its multi-scale anchor placement
anticipated the later feature pyramid approach, though the
features at different scales were extracted independently
without top-down refinement.

\subsubsection{YOLO Series}

Redmon \textit{et al.}~\cite{redmon2016you} reformulated
detection as a single regression from image pixels to box
coordinates and class probabilities, implemented as a single
CNN evaluated once per image.  YOLOv1 divided the image into a
$7 \times 7$ grid; each cell predicted $B$ bounding boxes and
$K$ class probabilities, yielding a fixed $S \times S \times
(B \cdot 5 + K)$ output tensor.  The key insight was that
removing the proposal stage allowed the entire detection system
to be trained end-to-end with a single loss function,
simplifying training and maximizing throughput.  YOLOv1
achieved 45 FPS at 63.4 mAP---substantially below Faster
R-CNN's accuracy, but establishing that extreme-throughput
detection was feasible.

YOLOv3~\cite{redmon2018yolov3} incorporated a Darknet-53
backbone, multi-scale prediction at three feature-map
resolutions, and anchor clustering, raising mAP to 33.0 on
COCO while maintaining real-time throughput.  YOLOv4
\cite{bochkovskiy2020yolov4} added a wide array of
bag-of-freebies and bag-of-specials training techniques---mosaic
augmentation, DropBlock, Mish activation, and a CSP backbone---
achieving 43.5 mAP at 65 FPS on COCO with a single RTX 2080 Ti,
a landmark result that substantially advanced the Pareto frontier.

\subsubsection{YOLOX and Decoupled Heads}

YOLOX~\cite{ge2021yolox} departed from the anchor-based
paradigm of earlier YOLO variants.  By removing anchor
definitions and adopting an anchor-free detection scheme,
YOLOX simplified the matching between ground-truth objects and
prediction locations.  Crucially, it introduced a
\textit{decoupled head}---separate $1\times1$ convolutional
sub-networks for classification and regression---resolving the
representational conflict that arises when a shared head is
asked to simultaneously optimize semantic discriminability and
spatial precision.  YOLOX-x achieved 51.2 mAP on COCO-val,
surpassing all prior YOLO variants while remaining real-time
capable at input resolution $640 \times 640$.

\subsubsection{YOLOv8}

The YOLOv8 architecture~\cite{wang2023yolov8} further refined
the single-stage paradigm with an updated C2f backbone module
(a cross-stage partial bottleneck with two convolutions),
a refined anchor-free head, and improved training recipes.
YOLOv8-x achieves 53.9 mAP on COCO-val at 14.3 ms per image,
establishing the immediate benchmark against which FusionDet
is compared in Section~\ref{sec:results}.

\subsection{Attention Mechanisms in Detection}

Attention mechanisms, originally introduced for sequence
modelling by the Transformer architecture~\cite{vaswani2017attention},
have been widely adapted for visual detection.  Squeeze-and-Excitation
(SE) networks~\cite{hu2018squeeze} model inter-channel
dependencies by globally pooling spatial information and
learning per-channel recalibration weights.  CBAM~\cite{woo2018cbam}
combines channel and spatial attention, sequentially
re-weighting feature maps along both dimensions.  Empirical
results across multiple detection architectures confirm that
incorporating attention into the feature neck---the module
responsible for aggregating features across scales---yields
consistent mAP improvements with minimal FPS penalty, a finding
that directly motivates the attention-augmented neck of
FusionDet.

\subsection{Feature Pyramid Networks}

Lin \textit{et al.}~\cite{lin2017feature} introduced the
Feature Pyramid Network (FPN), which added a top-down pathway
to the standard bottom-up CNN backbone.  Lateral connections
merge semantically rich, spatially coarse deep features with
spatially precise shallow features at each scale level, producing
a set of feature maps that are simultaneously informative and
well-localized.  FPN substantially improved small-object
detection and has since been incorporated into virtually every
competitive detection architecture.  PANet~\cite{liu2018path}
augmented FPN with an additional bottom-up pathway, shortening
the information path between low-level and high-level features.
The BiFPN variant used in EfficientDet introduced learnable
per-feature-level weighting, further improving multi-scale
fusion quality.  FusionDet builds upon these advances with a
bidirectional fusion neck incorporating attention recalibration,
as described in the following section.
